This literature review demonstrates how the dissertation emerges from the intersection of multiple theoretical traditions while offering novel methodological contributions:

\textbf{From Post-phenomenology and Agential Realism:} an understanding of technology as culturally mediated and materially embedded.
\textbf{From Speculative Design and Black Speculative Traditions:} methods for using imagination to open alternative technological relationships and futures.
\textbf{From Fugitive Theory and Black Radical Tradition:} strategies of refusal, appropriation, and transformation within constraining systems.
\textbf{From Embodied Knowledge and Performance Studies:} recognition of how bodies generate, transmit, and archive cultural knowledge.
\textbf{From Sonic Interaction and Multisensory Design:} insights into sound as a primary mode of interaction, perception, and embodied knowing.
\textbf{From Critical Technical Practice and Research-Creation:} approaches that integrate theory with hands-on experimentation and making.
\textbf{From Computational Culture Studies:} tools for analyzing how algorithms and code encode cultural assumptions.

The unique contribution lies in weaving these strands into a methodology that is both autoethnographic and collaborative: one that grounds technological invention in lived cultural practice and develops new procedural logics through small-scale, co-creative engagements. Rather than seeking universal frameworks, this approach privileges situated knowledge and collaborative invention as the means by which cultural epistemologies take computational form.